%!TEX root = ../nauchka.tex

\section{Планируемые эксперименты}

\begin{frame}
\frametitle{Цель экспериментов}

\begin{block}{Главный вопрос}
Подтвердить, что Tensor Parallelism решает проблему памяти и позволяет верифицировать модели, которые не влезают на один GPU
\end{block}

\vspace{0.5cm}

\textbf{Гипотеза:}
\begin{itemize}
    \item Модель размера $M$ вызывает OOM на 1 GPU
    \item Та же модель с TP=N успешно верифицируется на N GPU
    \item Модель размера $M \times N$ с TP=N использует примерно столько же памяти, как модель $M$ на 1 GPU
\end{itemize}

\end{frame}

\begin{frame}
\frametitle{Эксперимент 1: Базовый тест памяти}

\textbf{Задача:} Найти точку OOM для одного GPU

\vspace{0.3cm}

\textbf{Метод:}
\begin{enumerate}
    \item Взять простую модель (MLP или WideResNet)
    \item Постепенно увеличивать ширину слоев
    \item Для каждого размера запускать верификацию и измерять память
    \item Найти размер $M$, при котором происходит OOM
\end{enumerate}

\vspace{0.3cm}

\textbf{Ожидаемый результат:}
\begin{itemize}
    \item Определена граница $M$ для одного GPU
    \item Установлена baseline для сравнения с TP
\end{itemize}

\end{frame}

\begin{frame}
\frametitle{Эксперимент 2: TP масштабирование}

\textbf{Задача:} Проверить, что TP позволяет верифицировать модель размера $M \times N$ на N GPU

\vspace{0.3cm}

\textbf{Метод:}
\begin{enumerate}
    \item Взять модель размера $M$ (которая вызывает OOM на 1 GPU)
    \item Запустить с TP=2 на 2 GPU
    \item Проверить, что верификация проходит успешно
    \item Увеличить модель до размера $2M$ с TP=2
    \item Измерить потребление памяти на каждом GPU
\end{enumerate}

\vspace{0.3cm}

\textbf{Ожидаемый результат:}
\begin{itemize}
    \item Модель $M$ с TP=2 успешно верифицируется
    \item Модель $2M$ с TP=2 использует примерно столько же памяти на GPU, как модель $M$ на 1 GPU
    \item Доказано, что TP решает проблему масштабируемости
\end{itemize}

\end{frame}

\begin{frame}
\frametitle{Эксперимент 3: Точность границ}

\textbf{Задача:} Убедиться, что распределенное вычисление не влияет на точность границ

\vspace{0.3cm}

\textbf{Метод:}
\begin{enumerate}
    \item Взять модель, которая влезает на 1 GPU
    \item Вычислить границы обычным способом (без TP)
    \item Вычислить границы с TP=2
    \item Сравнить результаты (должны совпадать с точностью до floating point)
\end{enumerate}

\vspace{0.3cm}

\textbf{Ожидаемый результат:}
\begin{itemize}
    \item Границы совпадают с точностью $10^{-6}$
    \item Распределенное вычисление корректно
\end{itemize}

\end{frame}

\begin{frame}
\frametitle{Следующие шаги}

\begin{block}{Краткосрочные задачи}
\begin{enumerate}
    \item Завершить реализацию и тестирование на 2 GPU
    \item Провести эксперименты 1-3
    \item Интегрировать с полным циклом AlphaBetaCROWN
\end{enumerate}
\end{block}

\vspace{0.3cm}

\begin{block}{Долгосрочные задачи (к защите диплома)}
\begin{enumerate}
    \item Оптимизация коммуникаций (overlapping)
    \item Тестирование на реальных больших моделях (Transformers, LLM)
\end{enumerate}
\end{block}

\end{frame}